\clearpage\fancyhead[L]{\small VESTA.}
\begin{center}
\textbf{VESTA.}

\hrulefill

Monatliche Correspondenz. Band XV. S. 502–508. Mai 1807.

\hrulefill
\end{center}

--- --- ---

\par\indent „Sobald ich mit einiger Gewissheit annehmen konnte“, (schreibt uns Dr. \textsc{Olbers} am 22. April) „dass ich tatsächlich der erste Entdecker dieses neuen Planeten war, bat ich unseren unvergleichlichen \textsc{Gauss}, der sich um all diese kleinen Planeten so außerordentliche Verdienste erworben hat, einen Namen und ein Zeichen für ihn zu bestimmen. Dr. \textsc{Gauss} hat meine Bitte erfüllt. ‚Sie bitten mich‘, antwortete er, ‚so eindringlich, die Ehre anzunehmen und Pate Ihres Planeten zu sein, dass ich sie nicht ablehnen kann, so wenig Anspruch ich auch darauf habe. So sei es also. Ich weiß dem Planeten keinen schöneren Namen zu geben als den der Göttin, die die Völker der Antike zur Schutzgöttin der reinen Sitten, der makellosen Tugend und des häuslichen Glücks machten. Wenn Sie meine Wahl also nicht unpassend finden, dann heiße Ihr Töchterchen \emph{Vesta!}‘ Ich finde den Namen sehr glücklich gewählt. Als Zeichen hat Herr Dr. \textsc{Gauss} die symbolische Darstellung des heiligen Feuers gewählt, das auf dem Altar der Göttin brennt \VestaAltarBody{}. Auch dieses Zeichen scheint mir in jeder Hinsicht seinem Zweck zu entsprechen.“

\begin{flushright}
[\textsc{von Zach.}]
\end{flushright}
